\documentclass{article}

\textwidth=6.5in
\textheight=8.5in
\voffset=-54pt
\oddsidemargin=0in
\evensidemargin=0in
\setlength{\parskip}{8pt}
\setlength{\parindent}{0pt}

\usepackage{amsmath, amssymb}
\usepackage{ifthen}

\usepackage[inline]{enumitem}
\setlist{topsep=0pt, parsep=8pt}

\usepackage[rgb,table]{xcolor}\convertcolorsUfalse
\usepackage{graphicx}

% change hyperref options
\PassOptionsToPackage{hyphens}{url}
\usepackage[bookmarks,colorlinks,linkcolor=black,citecolor=blue,pdfstartview=FitH,urlcolor=blue]{hyperref}
\hypersetup{pdfpagemode=UseNone}
\usepackage{bookmark}

\usepackage{graphicx}
\usepackage{multicol}
\usepackage{framed}
\usepackage{array}

\usepackage{icomma}

\usepackage{pifont}

\usepackage{soul}

\usepackage{pgfplots}
\pgfplotsset{compat=1.18}  
\pgfplotscreateplotcyclelist{mycolor}{black, blue, red, green!70!black, magenta, purple, cyan, orange }  
\pgfplotsset{
  disabledatascaling,
  cycle list=tikzcolor, 
  axis lines=middle,
  every axis plot/.style={no marks, thick},
  every axis/.style={
    enlarge x limits=false,
    enlarge y limits=auto,
    xlabel style={at={(current axis.right of origin)},anchor=west},
    ylabel style={at={(current axis.above origin)},anchor=south},
    % extra x and y ticks for asymptotes
    extra x tick style={grid=major, grid style={dashed,black}},
    extra y tick style={grid=major, grid style={dashed,black}},
    /pgf/number format/1000 sep={},
  },    
  tick label style = {font=\footnotesize, fill=\currentbgcolor, inner sep=0pt},    
  xticklabel shift=1pt,    
  scaled ticks=false,
  scaled y ticks=false,
  unbounded coords=jump,
  samples=101,
  minor tick length=0,
  scatter/classes={
    open={mark=*,draw=black,fill=white, mark size=2, thin},
    closed={mark=*,draw=black,fill=black, mark size=2, thin}
  },
  width=3in,
}
\pgfplotsset{open/.style={only marks, mark=*, fill=white, thin, mark size=2}}
\pgfplotsset{closed/.style={only marks, mark=*, draw=black, fill=black,
    thin, mark size=2}}
\pgfplotsset{labeled/.style={nodes near coords, only marks, mark=*, draw=black, fill=black, thin, mark size=2, point meta=explicit symbolic}}

\usepackage{tikz}
\usetikzlibrary{
  matrix,%
  % enables the snake paths
  decorations.pathmorphing,%
  % enables bigger arrows
  decorations.markings,%
  decorations.pathreplacing,%
  decorations.text,%
  calc,%
  patterns,%
  shapes.geometric,%
  arrows,
  decorations.shapes,
  positioning,
  plotmarks,
  intersections
}
\tikzset{>=latex} % sets default arrow type in tikz
\tikzset{font=\normalsize}
\tikzset{mark size=1.5pt, mark options=thin}
\tikzset{pin distance=4pt,
  every pin edge/.style={<-, thin, shorten <= -2pt}}

\interfootnotelinepenalty=10000
\renewcommand{\thefootnote}{(\arabic{footnote})}

\usepackage{multirow, bigdelim}

% change to \usepackage[sols]{handout} to see solutions
\usepackage[sols, grading, workshop]{handout}

\providecommand{\check}{}
\renewcommand{\check}{\ifmmode{\text{ \ding{51} }}\else{\ding{51} }\fi}

\newcommand\iscurrentenumi[1]{%
  \ifthenelse{\equal{#1}{\theenumi}}{}{#1}}
\setenumerate[1]{ref=\#\arabic*}
\setenumerate[2]{ref=\protect\iscurrentenumi{\theenumi}(\alph*)}
\newcommand\enumiiprefix[2]{%
  \ifthenelse{{\equal{#1}{\theenumi}}\AND{\equal{#2}{(\alph{enumii})}}}{}{}%
  \ifthenelse{{\equal{#1}{\theenumi}}\AND{\NOT{\equal{#2}{(\alph{enumii})}}}}{#2}{}%
  \ifthenelse{\equal{#1}{\theenumi}}{}{#1#2}%
}
\setenumerate[3]{ref=\protect\enumiiprefix{\theenumi}{(\alph{enumii})}\roman*}

\makeatletter

% underline links               
\let\@href=\href
\renewcommand{\href}[2]{\@href{#1}{\ul{#2}}}

\makeatother

\definecolor{purple}{rgb}{0.5,0,1.0}

\newcommand{\doedfinity}[2][\newpage]{
  \begin{problem}[#1]
    Do the Edfinity assignment ``Problem Set #2''.

    We encourage you to write down any scratch work here, as well as
    things you want your future self to remember. Nothing you write
    here will be graded; it's just for you!
  \end{problem}
}

\newcommand{\psrnewpage}{\newpage}
\newcommand{\psreflection}[2][]{

  \ifthenelse{\not\equal{#1}{nocite}}{
    \begin{enumerate}[resume]
    \item
      \begin{problem}[1in]
        Please cite any help you got on this problem set:
        \begin{itemize}[nosep]
        \item If you worked with other students, please list their names here.
        \item If you used resources other than those on our Canvas site, please list them here.
        \end{itemize}
      \end{problem}

    \end{enumerate}
  }

  \probonly{%
    #2

    \psrnewpage

    \emph{Reflection space.} It's a good habit to take a few minutes at the end of each pset to reflect on what you learned and jot down notes to your future self. Here are a few reflection questions to get you started. Nothing you write here will be graded; it's just for you!
    \begin{itemize}
    \item Take a few minutes to look back at the problems. How could you explain to someone else how to do each problem? What was your strategy, and how did you know to use that strategy? If there are problems where you don't feel able to answer this, write down which ones they are. Come back to them in a few days when we post the homework solutions!

      \vspace{1.5in}

    \item Take a look back at the learning objectives on today's worksheet; how are you feeling about each one? If there are ones you know you need more practice with, write those down here.

      \vspace{1.5in}

    \item What did you find most challenging about this material? Do you have any lingering questions about it? If so, write them down here so you don't forget them. At some point, stop by office hours or MQC to ask!

      \vfill
    \end{itemize}      
  }
}

\usepackage{fancyhdr}
\lhead{\textcolor{white}{This content is protected and may not be shared, uploaded, or distributed.}}
\rhead{}
\rfoot{\vspace*{\fill}\footnotesize\textcopyright{} \emph{Harvard Math Ma}}
\renewcommand{\headrulewidth}{0pt}
\pagestyle{fancy}
\begin{document}

\begin{center}
  \large\textsc{Problem Set 13 - Limits}
\end{center}

\begin{enumerate}
\item  \note{
    \begin{itemize}[nosep]
    \item Limits from graphs
    \item Multiple choice question about size of $10^{30} \pm 10^{20}$ (to prepare for limits to infinity)
    \end{itemize}
}

  \doedfinity{13}

\item \label{something-over-0-limits}

  In class, you saw two kinds of limits that require extra care to compute. 

  \begin{enumerate}
  \item

    \begin{grading}
      1.5 points: 1 for correctly factoring and canceling, 0.5 for plugging in correctly to the simplified expression 
    \end{grading}

    \begin{problem}[\stretch{5}]
      The first kind is a limit of a fraction where the numerator and denominator both approach 0. In notation, we write that we're dealing with a limit of the form $\displaystyle \lim_{x \to a} \dfrac{f(x)}{g(x)}$ where $\displaystyle \lim_{x \to a} f(x) = 0$ and $\displaystyle \lim_{x \to a} g(x) = 0$.

      To evaluate such a limit, we try to simplify the fraction by factoring and canceling something in the numerator and denominator. Evaluate $\displaystyle \lim_{x \to 5} \frac{x^2 - 6x + 5}{x - 5}$ using this method. 
    \end{problem}

    \begin{solution}
      \begin{align*}
        \lim_{x \to 5} \frac{x^2 - 6x +5}{x-5} &= \lim_{x \to 5} \frac{(x-5)(x-1)}{x-5} \\
        &= \lim_{x \to 5} (x-1) \\
        &= 4
      \end{align*}
    \end{solution}

  \item The second kind is a limit of a fraction where the numerator goes to some non-zero number and the denominator goes to 0. In notation, we write that this is a limit of the form $\displaystyle \lim_{x \to a} \dfrac{f(x)}{g(x)}$ where $\displaystyle \lim_{x \to a} f(x) = L \neq 0$ and $\displaystyle \lim_{x \to a} g(x) = 0$.

    The issue with such a limit is that, when we divide a number close to $L$ by a number close to 0, we could get either a really big positive number or an extremely negative number. To deal with this, we should look at both one-sided limits. For each side, we want to think about whether the denominator is positive or negative. Let's do this for $\displaystyle \lim_{x \to 2} \frac{x + 3}{x - 2}$.

    \begin{enumerate}
    \item

      \begin{grading}
        2 points: 0.5 for the first question, 1 for the second, 0.5 for the third. 
      \end{grading}

      \begin{problem}[0.25in]
        Let's first look at $\displaystyle \lim_{x \to 2^+} \frac{x + 3}{x - 2}$. If $x$ is slightly bigger than 2, is $x - 2$ positive or negative? 
      \end{problem}

      \begin{solution}
        When $x$ is slightly bigger than $x$, $x-2$ is positive.
      \end{solution}

      \begin{problem}[\stretch{2}]
        Based on that, when $x$ is slightly bigger than 2, is $\dfrac{x + 3}{x - 2}$ a really big positive number or a very negative number? (Be sure to also describe the sign of the numerator when $x$ is slightly bigger than 2.) 
      \end{problem}

      \begin{solution}
        When $x$ is slightly bigger than $2$, $\dfrac{x+3}{x-2}$ is a positive number divided by a small
        positive number, which is a large positive number.
      \end{solution}

      \begin{problem}[\newpage]
        Based on that, is $\displaystyle \lim_{x \to 2^+} \frac{x + 3}{x - 2}$ equal to $\infty$ or $-\infty$? 
      \end{problem}

      \begin{solution}
        Based on that, $\displaystyle \lim_{x \to 2^+} \frac{x + 3}{x - 2}=\infty.$
      \end{solution}

    \item

      \begin{grading}
        1 point
      \end{grading}

      \begin{problem}[\stretch{2}]
        Use the same line of reasoning to find $\displaystyle \lim_{x \to 2^-} \frac{x + 3}{x - 2}$. Write out your reasoning.
      \end{problem}

      \begin{solution}
        When $x$ approaches $2$ from the left, the denominator $x-2$ takes on negative values close to zero,
        and the numerator is positive. This means that $\dfrac{x+2}{x-2}$ is taking on lower
        and lower negative values, decreasing without bound. Therefore 
        $\displaystyle \lim_{x \to 2^-} \frac{x + 3}{x - 2} = -\infty.$
      \end{solution}

    \item

      \begin{grading}
        0.5 points
      \end{grading}

      \begin{problem}[\stretch{0.75}]
        Now that you know both one-sided limits, what's $\displaystyle \lim_{x \to 2} \frac{x + 3}{x - 2}$? 
      \end{problem}

      \begin{solution}
        As $x$ approaches $2$, $\dfrac{x+3}{x-2}$ increases without bound on one side and decreases without
        bound on the other, so $\displaystyle \lim_{x \to 2} \frac{x + 3}{x - 2}$ does not exist.
      \end{solution}

    \end{enumerate}

  \end{enumerate}

\item \label{evaluate-limits-1} Evaluate the following limits without using technology. Be alert for the two kinds of limits you looked at in \ref{something-over-0-limits}! As usual, show your work or explain your reasoning in each part. 

  \begin{enumerate}

  \item \label{evaluate-limits-1-0/0}

    \begin{grading}
      2 points: 0.5 for realizing they need to factor and cancel, 1 for doing that correctly, 0.5 for plugging in correctly to the simplified expression 
    \end{grading}

    \begin{problem}[\vfill]
      $\displaystyle \lim_{x \to 3} \frac{x - 3}{x^2 - 5x + 6}$
    \end{problem}

    \begin{solution}
      When we try direct substitution, we get the undefined expression $\frac{0}{0}$. This tells us to try
      to factor and cancel.
      \begin{align*}
        \lim_{x \to 3} \frac{x-3}{x^2-5x+6} &= \lim_{x \to 3} \frac{x-3}{(x-3)(x-2)}\\
        &= \lim_{x \to 3} \frac{1}{x-2} \\
        &= 1
      \end{align*}
    \end{solution}

  \item

    \begin{grading}
      1 point
    \end{grading}

    \begin{problem}[\vfill]
      $\displaystyle \lim_{x \to 1} \frac{x - 3}{x^2 - 5x + 6}$
    \end{problem}

    \begin{solution}
      When we try direct substitution, we get $\frac{-2}{2}=-1$, so 
      $\displaystyle \lim_{x \to 1} \frac{x - 3}{x^2 - 5x + 6}=-1$.
    \end{solution}

  \item

    \begin{grading}
      1 point
    \end{grading}

    \begin{problem}[\vfill]
      $\displaystyle \lim_{x \to 10} \frac{x - 10}{(x - 4)^2}$
    \end{problem}

    \begin{solution}
      Direct substitution gives us $\frac{0}{36}=0$, so 
      $\displaystyle \lim_{x \to 10} \frac{x - 10}{(x - 4)^2} = 0.$
    \end{solution}

  \item \label{evaluate-limits-1-nonzero/0-1}

    \begin{grading}
      2 points. If they give some correct explanation but don't say anything about the sign of $(x - 4)^2$, deduct 1.
    \end{grading}

    \begin{problem}[\nstretch{1.5}]
      $\displaystyle \lim_{x \to 4} \frac{x - 10}{(x - 4)^2}$
    \end{problem}

    \begin{solution}
      Direct substitution gives us $\frac{-6}{0}$, which tells us that we should consider whether the function is 
      increasing or decreasing without bound. As $x$ approaches $4$, regardless of the direction,
      the numerator is negative, and the denominator
      is taking on positive values closer and closer to zero.
      As $x$ approaches $4$, $\dfrac{x - 10}{(x - 4)^2}$ is decreasing without bound, 
      so $\displaystyle \lim_{x \to 4} \frac{x - 10}{(x - 4)^2}=-\infty.$
    \end{solution}

  \item

    \begin{grading}
      3 points: 0.5 for recognizing they need to look at the sign of the denominator, 1 for reasoning correctly about the right-handed limit, 1 for reasoning correctly about the left-handed limit, 0.5 for a conclusion that matches their one-sided limits.

      If they just say $\infty$ without any of the above, give them 0.5 out of 3. 
    \end{grading}

    \begin{problem}[\vfill]
      $\displaystyle \lim_{x \to 3} \frac{5}{x^2 - 9}$
    \end{problem}

    \begin{solution}
      Direct substitution gives us $\frac{5}{0}$, which tells us that we should consider whether the function
      is increasing or decreasing without bound. As $x$ approaches $3$, the numerator is positive. When $x$ is
      less than $3$, $x^2-9$ is negative, so $\dfrac{5}{x^2-9}$ will be negative, thus
      $\displaystyle \lim_{x \to 3^-} \frac{5}{x^2-9} = -\infty$. When $x$ is greater than $3$, $x^2-9$ will be
      positive, so $\dfrac{5}{x^2-9}$ will be positive, 
      thus $\displaystyle \lim_{x \to 3^-} \frac{5}{x^2-9} = \infty$.
      This means $\displaystyle \lim_{x \to 3} \frac{5}{x^2-9}$ does not exist.
    \end{solution}

  \end{enumerate}

  \begin{problemonly}
    \check We suggest that you graph the functions $\dfrac{x - 3}{x^2 - 5x + 6}$, $\dfrac{x - 10}{(x - 4)^2}$, and $\dfrac{5}{x^2 - 9}$ using a tool like \href{https://www.desmos.com/calculator}{Desmos} to check your answers. (But be sure your justification doesn't rely on the graphs.) 
\end{problemonly}

\item All of the statements below are false. In each part, sketch the graph of a counterexample (that is, an example showing the statement is false); be sure to label relevant values on the axes.

  \begin{grading} 
    2 points per part: 1 for a sketch satisfying the ``if'' part of the if-then statement and 1 for the sketch not satisfying the ``then'' part.

    Of course, their examples may look very different from the ones in the solutions. 
  \end{grading} 

  \begin{enumerate}
  \item % 7.2.16

    \begin{problem}[\vfill]
      If $\displaystyle \lim_{x \rightarrow 5} f(x) = 2$, then $f(5)=2$.
    \end{problem}

    \begin{solution} 

      \begin{center}
        \begin{tikzpicture} 
          \begin{axis}[xtick={5}, ytick={2}]
            \addplot+[domain=-2:7] {(.04*x^3 -3};
            \addplot[closed] coordinates{(5,8)};
            \addplot[open] coordinates{(5,2)};
          \end{axis}
        \end{tikzpicture}
      \end{center}    

    \end{solution} 

  \item
    \begin{problem}[\vfill\newpage]
      If $\displaystyle \lim_{x \to 4} f(x) = 3$, then $f(4) \neq 3$. 
    \end{problem}

    \begin{solution}

      \begin{center}
        \begin{tikzpicture}
          \begin{axis}[xtick={4}, ytick={3}]
            \addplot+[domain=-2:7]{x - 1}; 
          \end{axis}
        \end{tikzpicture}
      \end{center}
    \end{solution}

  \item % 7.2.14

    \begin{problem}[\vfill]
      If $f(2)$ is undefined, then $\displaystyle \lim_{x \rightarrow 2} f(x)$ is undefined.
    \end{problem}

    \begin{solution} 

      \begin{center}
        \begin{tikzpicture} 
          \begin{axis}
            \addplot+[domain=-4:4] {(x^2};
            \addplot[open] coordinates{(2,4)};

          \end{axis}
        \end{tikzpicture}
      \end{center}    

    \end{solution}

  \item % 7.2.15
    \begin{problem}[\vfill]
      If $\displaystyle\lim_{x \rightarrow 0^+} f(x)$ and $\displaystyle \lim_{x \rightarrow 0^-} f(x)$ both exist, then $\displaystyle \lim_{x \rightarrow 0} f(x)$ exists.
    \end{problem}

    \begin{solution} 
      \begin{center}
        \begin{tikzpicture} 
          \begin{axis}
            \addplot+[domain=0:2] {(x+7};
            \addplot+[domain=-2:0] {(x^2};
            \addplot[open] coordinates{(0,0)};
            \addplot[closed] coordinates{(0,7)};
          \end{axis}
        \end{tikzpicture}
      \end{center}    
    \end{solution} 

  \end{enumerate} 

\end{enumerate}
\psreflection{}

\end{document}
